\section{An alternative construction using auxiliary double covers}
\label{app:branched}

We give an alternative construction of a manifold satisfying
Proposition~\ref{prop:transfer}.  Throughout this appendix, $X$ is a
connected smooth projective complex fourfold, $D\subset X$ is a nonempty
smooth connected reduced divisor, and $\mathcal X$ is its order-two root
orbifold. An affine action on a product of elliptic curves makes the
involution associated with $D$ act freely while retaining selected auxiliary fixed
loci. Blowing up these loci produces a smooth quotient with the ordinary
submersions needed along the auxiliary branch divisors.

\begin{lemma}[An auxiliary branched cover]
\label{lem:branched-cover}
There are distinct smooth prime divisors $B_1,\ldots,B_5$, with
$D+\sum_i B_i$ simple normal crossing, and a connected smooth projective
Galois cover
\[
 \pi:S\longrightarrow X,\qquad
 A=(\mathbb Z/2\mathbb Z)^5=\langle e_1,\ldots,e_5\rangle,
\]
with the following local descriptions.  Put $\tau=e_1+\cdots+e_5$ and
$I(x)=\{i:x\in B_i\}$.
\begin{enumerate}
\item If $x\notin D$, a chart over $x$ takes independent square roots of
the coordinates defining $B_i$, $i\in I(x)$.  Its stabilizer is generated
by the corresponding $e_i$, each changing the sign of its own root
coordinate.
\item If $x\in D$, choose coordinates $z_0,z_i$ defining $D,B_i$ for
$i\in I(x)$.  The charts over $x$ take independent roots
\[
 w_0^2=z_0,\qquad w_i^2=z_i\quad(i\in I(x)),
\]
and leave the remaining coordinates unchanged.  Their stabilizer is
$\langle\tau,e_i:i\in I(x)\rangle$; here $\tau$ changes the sign of
$w_0$ alone, and $e_i$ changes the sign of $w_i$ alone.
\end{enumerate}
In particular, $\pi^*D$ is twice a reduced divisor, and
$S^{e_i}$ is precisely the reduced inverse image of $B_i$, a smooth
divisor all of whose components dominate $B_i$.
\end{lemma}

\begin{proof}
Choose a line bundle $H$ such that $H^{\otimes2}(-D)$ is very ample.
Successive general members $B_i$ of this system are smooth and
irreducible and can be chosen transverse to all the preceding strata,
including those in $D$.  Thus the asserted simple normal crossing
condition follows from Bertini's theorem.  In particular, at a point of
$D$ at most three of the $B_i$ pass through the point.

Let $s_D$ and $s_i$ be the divisor sections for $D$ and $B_i$.  Form the
cover by adjoining roots
\[
 y_i^2=s_Ds_i\quad(1\leq i\leq5),
\]
where $y_i$ takes values in $H$, and normalize.  In a rational frame of
$H$, the five right-hand sides have independent classes modulo squares
in $\mathbb C(X)^*$.  Indeed, the order along $B_i$ of a product of these
functions is odd exactly when its $i$th factor occurs an odd number of
times.  A product representing a square must therefore have every
exponent even.  The resulting field extension has degree $2^5$, with
the independent sign changes as its Galois group.  Its normalization
$S$ is connected and finite over $X$, hence projective.

We compute this normalization locally analytically.  Units in the
equations have holomorphic square roots after shrinking the
neighborhood.  Off $D$, the equations give independent roots of the
passing $B_i$ coordinates; the other roots separate the local sheets.
At $x\in D$, the equations become
\[
 y_i^2=z_0z_i\quad(i\in I(x)),\qquad
 y_j^2=z_0\quad(j\notin I(x)).
\]
There is an index $j_0\notin I(x)$.  On a component set $w_0=y_{j_0}$
and, in its meromorphic function field, set $w_i=y_i/w_0$ for
$i\in I(x)$.  These functions satisfy the displayed independent root
equations.  The other $y_j$ equal $\pm w_0$, and their relative signs
separate the components.  The independent-root charts are finite and
normal over the base chart and have exactly the prescribed generic
sheets, so they are its normalization.  They are smooth.  The sign
changes give the stabilizers and their actions stated in the lemma.

For completeness, every element of the latter stabilizer with
nonzero $\tau$-coefficient has at least
\[
 5-|I(x)|\geq2
\]
nonzero coordinates in the basis $e_1,\ldots,e_5$.  Consequently $e_i$
belongs to a point stabilizer precisely over $B_i$.  In the local
charts its fixed locus is the root-coordinate hyperplane for $B_i$.
This proves the assertion about $S^{e_i}$.  Each of its components has
dimension three, so its finite image in the irreducible divisor $B_i$
has dimension three and equals $B_i$.  The equation $z_0=w_0^2$ also
proves the assertion about $\pi^*D$.
\end{proof}

\begin{lemma}[An action on a product of elliptic curves]
\label{lem:branched-action}
There is an action of $A$ on an abelian variety $F$ of dimension ten
such that every element with at least two nonzero coordinates in the
basis $e_i$ is fixed-point-free.  Each $e_i$ has a nonempty smooth fixed
locus of codimension four, and these five fixed loci are pairwise
disjoint.
\end{lemma}

\begin{proof}
Fix an elliptic curve $E$ and a nonzero point $c\in E[2]$.  Set
\[
 F=\prod_{1\leq i<j\leq5}E_{ij},\qquad E_{ij}=E.
\]
On the factor $E_{ij}$ let $e_i$ act by $v\mapsto-v$, let $e_j$ act by
$v\mapsto c-v$, and let the other generators act trivially.  The two
displayed involutions commute because $2c=0$, and their product is
translation by $c$.  If an element has both $i$ and $j$ in its support,
it therefore acts without fixed points on this factor and hence on
$F$.

For a single $e_i$, its four moving factors each carry an involution
of the form $v\mapsto a-v$.  The equation $2v=a$ has solutions and a
finite reduced solution set.  Thus $F^{e_i}$ is a nonempty disjoint
union of translates of products of the six remaining elliptic
factors.  It is smooth of codimension four.  Two different generators
cannot fix the same point, since their product is fixed-point-free.
\end{proof}

\begin{proposition}[A manifold model of dimension fourteen]
\label{prop:branched-model}
There is a connected smooth projective variety $Y$ of dimension
fourteen and a surjective morphism $f:Y\to X$ with the following
properties.
\begin{enumerate}
\item Over
\[
 U=X\setminus\bigl(D\cup B_1\cup\cdots\cup B_5\bigr),
\]
the map is a holomorphically locally trivial proper bundle with
connected fibre $F$, the abelian variety of
Lemma~\ref{lem:branched-action}.
\item The divisor $f^*D$ is twice a reduced divisor.  At every point
over $D$, the map admits a local holomorphic lift to its root chart.
At points over a general point of $D$, such a lift can be chosen
submersive.
\item Above a general point of every prime divisor of $X$ other than
$D$, there is a point at which $f$ is a submersion to $X$.
\end{enumerate}
\end{proposition}

\begin{proof}
Use the diagonal action of $A$ on $V=S\times F$.  By
Lemma~\ref{lem:branched-action}, a nontrivial point stabilizer can only
be generated by a single $e_i$.  The entire nonfree locus is the union
of the pairwise disjoint smooth subvarieties
\[
 Z_i=S^{e_i}\times F^{e_i}\qquad(1\leq i\leq5).
\]
They have codimension five.  They are invariant under $A$, since $A$
is abelian.  At a point of $Z_i$ the stabilizer is exactly
$\langle e_i\rangle$.  In local coordinates it fixes the tangent
coordinates to $Z_i$ and negates all five normal coordinates: one
from $S$, and four from the moving elliptic factors.

Blow up their disjoint union and take the quotient:
\[
 b:\widehat V=\operatorname{Bl}_{\cup_i Z_i}(V)\longrightarrow V,
 \qquad q:\widehat V\longrightarrow Y=\widehat V/A.
\]
The action lifts to the blowup.  We verify that $Y$ is smooth.  If
$(v_0,\ldots,v_4)$ are the normal coordinates just described, then in
the blowup chart with distinguished coordinate $v_j$ we have
coordinates
\[
 t=v_j,\qquad a_k=v_k/v_j\quad(k\ne j),
\]
besides the tangent coordinates.  The stabilizer acts by
\[
 (t,(a_k),\text{tangent coordinates})
 \longmapsto
 (-t,(a_k),\text{tangent coordinates}).
\]
Its quotient chart has smooth coordinates $t^2,(a_k)$ and the tangent
coordinates.  Every stabilizer on $\widehat V$ is contained in the
stabilizer of its image in $V$; away from the exceptional divisors it
is trivial.  These charts therefore establish smoothness everywhere.

The blowup is connected, smooth, and projective.  Its finite-group
quotient is projective: tensoring the translates of an ample line
bundle gives an invariant linearized ample line bundle, whose
invariant section ring constructs the projective quotient.  The
quotient is connected and has the same dimension as $V$, namely
$4+10=14$.  The morphism
$\pi\circ\operatorname{pr}_S\circ b$ is invariant and descends to a
surjective morphism $f:Y\to X$.

Over $U$, the map $S\to X$ is an \'etale $A$-torsor and the blowup has
no centre.  Hence
\[
 f^{-1}(U)=(S|_U\times F)/A.
\]
Analytically trivializing the finite \'etale torsor identifies this
map locally with projection from a product with $F$.  This proves the
first assertion, including properness, smoothness, and connectedness
of these fibres.  Their cotangent bundles are trivial.

We next check the divisor $D$ at all points, not only at generic
points away from the auxiliary divisors.  Let $x\in D$ and use the
coordinates of Lemma~\ref{lem:branched-cover}.  A stabilizer at a point
of $S\times F$ over $x$ is either trivial or $\langle e_i\rangle$ for
some $i\in I(x)$.  Indeed an inertia element involving $\tau$ has
support of size at least two and has no fixed point on $F$.  In
particular, every possible product stabilizer fixes the function
$w_0$, whose square is the pullback of $z_0$.

The pullback of $w_0$ to a blowup chart is still holomorphic and is
fixed by the stabilizer of every point in that chart.  A neighborhood
of the corresponding point of the finite quotient is the quotient
of a sufficiently small neighborhood by this stabilizer.  Thus
$w_0$ descends there to a holomorphic function whose square is
$f^*z_0$.  Together with the unchanged base coordinates it gives a
local holomorphic lift to the root chart of $D$.  This argument
includes points on exceptional divisors and all intersections with
the auxiliary branch divisors.

We can also determine the multiplicities exactly.  Write
$\pi^*D=2R$ with $R$ reduced.  No component of a centre $Z_i$ is
contained in $R\times F$, because every component of $S^{e_i}$
dominates $B_i\ne D$.  Hence blowing up introduces no exceptional
component in the pullback of $R\times F$, and
\[
 (f q)^*D=2\widehat R,
\]
where $\widehat R$ is its reduced strict transform.  At a general
point of every component of $\widehat R$ the image in $D$ avoids all
$B_i$.  There the only nontrivial inertia of $S\to X$ is generated
by $\tau$, which acts freely on $F$, so $q$ is unramified.  Every
prime component of $f^*D$ is dominated by a prime component of
$\widehat R$, since $q$ is finite.  The ramification index there is
one, and the displayed equality proves that its coefficient in
$f^*D$ is exactly two.

Over a general point of $D$, all the auxiliary divisors are absent.
There are no blowup centres and the diagonal action is free.
Consequently the quotient is locally isomorphic to $S\times F$, and
projection to $(w_0,z_2,z_3,z_4)$ is a submersive local lift to the
root chart.  This completes the second assertion.

Finally fix $i$ and take a general point of $B_i$ outside $D$ and the
other $B_j$.  Choose a point of $S^{e_i}$ above it and a point of
$F^{e_i}$.  Near the resulting point of $Z_i$, let $w$ be the cover
coordinate with $w^2=z_i$, where $z_i$ defines $B_i$.  The five normal
coordinates to the centre are $w,u_1,\ldots,u_4$, with the $u_j$
coming from the moving elliptic factors.  The other three base
coordinates are tangent to the centre and remain unchanged.

In the blowup chart distinguished by $w$, use
\[
 w,\quad u_1/w,\ldots,u_4/w
\]
as its normal-direction coordinates.  Passing to the local quotient
replaces $w$ by $\xi=w^2$ and leaves the ratios and all tangent
coordinates unchanged.  The map to $X$ then has coordinates
\[
 z_i=\xi,
 \qquad\text{the other three base coordinates unchanged}.
\]
It is a submersion, including at $\xi=0$ on the exceptional divisor.
This provides the required point above each general point of $B_i$.
Every other prime divisor distinct from $D$ has its general point in
$U$, where $f$ is smooth.  The third assertion follows.
\end{proof}

If $\mathcal X$ is special, Proposition~\ref{prop:transfer} now shows
that this $Y$ is special and gives an epimorphism
\[
 \pi_1(Y)\longrightarrow
 \pi_1(X\setminus D)/\langle\!\langle\gamma^2\rangle\!\rangle.
\]
Assumption~\ref{ass:manifold}, applied to this special smooth projective
fourteenfold, makes $\pi_1(Y)$ virtually abelian. Its finite-index
abelian subgroup has finite-index abelian image under the displayed
surjection. This gives a second proof of Theorem~\ref{thm:main}.
